% ENTER27 Conference Presentation extended abstract.
% Limit: 2 pages including key references. Springer LNCS class (llncs v2.25).
% Review version is double-blind: do not restore author names or affiliations.
% This draft reports the design and the frozen screening status. It does not
% report creator differences. The formal analysis sample is not yet frozen.
\documentclass[runningheads]{llncs}
\usepackage[T1]{fontenc}
\usepackage{newtxtext}
\usepackage[varvw]{newtxmath}

\begin{document}

\title{Who Represents the Destination? Separating Creator Evidence from What Tourism Posts Show}
\titlerunning{Creator Evidence and Tourism Posts}

\author{Anonymous Author}
\authorrunning{Anonymous Author}
\institute{Author names and affiliations are omitted for double-blind review.}

\maketitle

\begin{abstract}
Dashboards and classifiers now treat social-media posts as if they were the destination image held by tourists. This extended abstract sets out a production-side design that keeps the two apart. In a single-destination study of Qingdao, it asks how creator groups differ in the strategies they use and in the resources, attributes, and evaluations those strategies organize. Account subject, public positioning, and audience size are recorded separately from the content codes, and the comparison groups will be locked before any group result is inspected. A five-platform archive has been screened, but human adjudication is still open, so no creator difference is reported here.

\keywords{user-generated content \and destination image \and content analysis \and creator types \and social media.}
\end{abstract}

\section{Background}
Destination image is the set of beliefs, ideas, and impressions a person holds about a place~\cite{crompton1979}. Attribute-based and holistic measures do not collapse into one content count~\cite{echtner1993,tasci2007}. Posts supply material for image formation~\cite{gartner1994} and can be coded for cognitive and affective attributes~\cite{garay2019}, but a post is not the image a visitor holds, and it does not show persuasion or spending. Classifiers and dashboards often skip that boundary and treat the producing account as a generic consumer. A destination organization then confuses content volume with evidence of belief or of campaign return.

\section{Purpose}
The study asks how different types of tourism-content creators use observable strategies to represent Qingdao. Three questions follow. Which strategies do the groups use? Which resources and evaluated attributes do they bring forward, and with what direction and strength? How do strategies co-occur with resources and evaluations inside the same textual segment? Co-occurrence is a property of the finished post. It is not evidence of intent, and it is not an effect on audiences. Subject, positioning, and follower count are measured separately so that the grouping rule can be locked before the content results are seen.

\section{Method}
The method is theory-led structured content analysis~\cite{krippendorff2018,neuendorf2017}. One city gives creators a shared resource base. Text is the primary evidence. Posts are split into stable segments, and each segment receives one dominant evaluation. Coders record the resources mentioned, the mode of presentation, the direction and strength of evaluation, and the attribute judged. Two coders work independently; adjudication is stored apart from the first answers. Images enter only after their own quality checks. A later model may apply accepted human labels to uncoded posts. It does not define the categories.

On 12~September 2026 the project froze an archive of 19{,}984 posts and 161{,}377 image links from Bilibili, Douyin, Weibo, Xiaohongshu, and Zhihu. Of these, 19{,}255 posts passed a topic screen. A frozen classifier then marked 10{,}914 for retention, 5{,}305 for exclusion, and 3{,}036 for manual review. Every candidate is still awaiting human adjudication. The formal analysis sample is empty. The counts describe the archive and the screening pool. They are not a result about creators.

\section{Significance}
The claim comes before any group finding. Analytics that price content volume, or that treat a generated label as the visitor's image, measure the wrong object. Separating who posted from what the post shows tells a destination organization which evidence can support a decision. Next come coder agreement, a grouping rule locked with enough accounts in each group, and only then the comparisons. Until those gates are passed, the study will not report a creator difference and will not scale labels by model.

\begin{thebibliography}{7}
\bibitem{crompton1979}
Crompton, J.L.: An assessment of the image of Mexico as a vacation destination. J. Travel Res. \textbf{18}(4), 18--23 (1979)

\bibitem{echtner1993}
Echtner, C.M., Ritchie, J.R.B.: The measurement of destination image: An empirical assessment. J. Travel Res. \textbf{31}(4), 3--13 (1993)

\bibitem{garay2019}
Garay, L.: \#Visitspain. Breaking down affective and cognitive attributes in the social media construction of the tourist destination image. Tour. Manag. Perspect. \textbf{32}, 100560 (2019)

\bibitem{gartner1994}
Gartner, W.C.: Image formation process. J. Travel Tour. Mark. \textbf{2}(2--3), 191--216 (1994)

\bibitem{krippendorff2018}
Krippendorff, K.: Content Analysis: An Introduction to Its Methodology, 4th edn. Sage, Thousand Oaks (2018)

\bibitem{neuendorf2017}
Neuendorf, K.A.: The Content Analysis Guidebook, 2nd edn. Sage, Thousand Oaks (2017)

\bibitem{tasci2007}
Tasci, A.D.A., Gartner, W.C., Cavusgil, S.T.: Conceptualization and operationalization of destination image. J. Hosp. Tour. Res. \textbf{31}(2), 194--223 (2007)
\end{thebibliography}

\end{document}
